% GENERATED FILE. Edit canonical lessons, then run npm run generate:books.
% canonical-source-hash: fnv1a64:f0c91dcfea03f7af
% canonical-lessons: GE-C520-krimi, GE-C520-biografie, GE-C520-kurzgeschichte, GE-C520-sachbuch, GE-C520-boesewicht

\chapter{Krimis und Biografien}
\label{ch:ge-krimis-und-biografien}
\hlchaptermodality{520}

\begin{chapteropening}
\textbf{By the end of this chapter:} \emph{I can talk about kinds of books.}

Complete the last lesson of chapter 520: I can talk about kinds of books.
\end{chapteropening}

\section[der Krimi]{der Krimi — crime novel, crime film}

\label{lesson:GE-C520-krimi}



\begin{quote}
Before the new one: say the German for sculptor, then the German for work of art.
\end{quote}

\subsection*{You'll want to know: der Krimi}
\textbf{der Krimi} — \textquotedblleft{}crime novel, crime film\textquotedblright{}.

\begin{quote}
Am Strand lese ich gern einen spannenden Krimi.
\end{quote}

\begin{quote}
\emph{On the beach I like reading an exciting crime novel.}
\end{quote}

Plural: \emph{die Krimis}. Short for \emph{der Kriminalroman}.

\begin{cousinweb}
\emph{Krimi} is short for \emph{Kriminalroman} \textquotedblleft{}crime novel\textquotedblright{}, built on a learned word from Latin \emph{crīmen} \textquotedblleft{}charge, crime\textquotedblright{} (Lewis \& Short).
\end{cousinweb}

\begin{grammarlens}[title={What you've built}]
One more word for telling, explaining and describing at B1.
\end{grammarlens}

\subsection*{Your turn}
\begin{itemize}
\raggedright
  \item \emph{Say it:} \emph{der Krimi}
  \item \emph{Say it:} \emph{der Krimi}, once more
  \item \emph{Say it:} a sentence of your own with \emph{der Krimi}
  \item \emph{Recall:} say the German for sculptor, then the German for work of art, then say \emph{der Krimi} again
\end{itemize}

\begin{tcolorbox}[breakable,colback=teal!4,colframe=teal!35!black,title={Before you move on}]
What does der Krimi mean? (\textquotedblleft{}Crime novel, crime film\textquotedblright{}.) Say it once more.
\end{tcolorbox}
\section[die Biografie]{die Biografie — biography}

\label{lesson:GE-C520-biografie}



\begin{quote}
Before the new one: say the German for work of art, then the German for crime novel.
\end{quote}

\subsection*{You'll want to know: die Biografie}
\textbf{die Biografie} — \textquotedblleft{}biography\textquotedblright{}.

\begin{quote}
Sie liest eine Biografie über eine berühmte Sängerin.
\end{quote}

\begin{quote}
\emph{She is reading a biography of a famous singer.}
\end{quote}

Plural: \emph{die Biografien}. Also spelt \emph{die Biographie}.

\begin{cousinweb}
\emph{Biografie} is a learned word built from Greek \emph{bíos} \textquotedblleft{}life\textquotedblright{} and \emph{gráphein} \textquotedblleft{}to write\textquotedblright{}; EtymDB lists the German word without a source.
\end{cousinweb}

\begin{grammarlens}[title={What you've built}]
One more word for telling, explaining and describing at B1.
\end{grammarlens}

\subsection*{Your turn}
\begin{itemize}
\raggedright
  \item \emph{Say it:} \emph{die Biografie}
  \item \emph{Say it:} \emph{die Biografie}, once more
  \item \emph{Say it:} a sentence of your own with \emph{die Biografie}
  \item \emph{Recall:} say the German for work of art, then the German for crime novel, then say \emph{die Biografie} again
\end{itemize}

\begin{tcolorbox}[breakable,colback=teal!4,colframe=teal!35!black,title={Before you move on}]
What does die Biografie mean? (\textquotedblleft{}Biography\textquotedblright{}.) Say it once more.
\end{tcolorbox}
\section[die Kurzgeschichte]{die Kurzgeschichte — short story}

\label{lesson:GE-C520-kurzgeschichte}



\begin{quote}
Before the new one: say the German for crime novel, then the German for biography.
\end{quote}

\subsection*{You'll want to know: die Kurzgeschichte}
\textbf{die Kurzgeschichte} — \textquotedblleft{}short story\textquotedblright{}.

\begin{quote}
In der Schule lesen wir eine Kurzgeschichte.
\end{quote}

\begin{quote}
\emph{At school we are reading a short story.}
\end{quote}

Plural: \emph{die Kurzgeschichten}.

\begin{cousinweb}
\emph{Kurzgeschichte} is built on \emph{kurz} \textquotedblleft{}short\textquotedblright{}, an early borrowing from Latin \emph{curtus}, and \emph{Geschichte} \textquotedblleft{}story\textquotedblright{}, from Middle High German \emph{geschicht} \textquotedblleft{}event, story\textquotedblright{} (EtymDB).
\end{cousinweb}

\begin{grammarlens}[title={What you've built}]
One more word for telling, explaining and describing at B1.
\end{grammarlens}

\subsection*{Your turn}
\begin{itemize}
\raggedright
  \item \emph{Say it:} \emph{die Kurzgeschichte}
  \item \emph{Say it:} \emph{die Kurzgeschichte}, once more
  \item \emph{Say it:} a sentence of your own with \emph{die Kurzgeschichte}
  \item \emph{Recall:} say the German for crime novel, then the German for biography, then say \emph{die Kurzgeschichte} again
\end{itemize}

\begin{tcolorbox}[breakable,colback=teal!4,colframe=teal!35!black,title={Before you move on}]
What does die Kurzgeschichte mean? (\textquotedblleft{}Short story\textquotedblright{}.) Say it once more.
\end{tcolorbox}
\section[das Sachbuch]{das Sachbuch — non-fiction book}

\label{lesson:GE-C520-sachbuch}



\begin{quote}
Before the new one: say the German for biography, then the German for short story.
\end{quote}

\subsection*{You'll want to know: das Sachbuch}
\textbf{das Sachbuch} — \textquotedblleft{}non-fiction book\textquotedblright{}.

\begin{quote}
Zum Geburtstag bekam ich ein Sachbuch über Vögel.
\end{quote}

\begin{quote}
\emph{For my birthday I got a non-fiction book about birds.}
\end{quote}

Plural: \emph{die Sachbücher}. Fiction is \emph{die Belletristik}.

\begin{cousinweb}
\emph{Sachbuch} is built on \emph{Sache} \textquotedblleft{}thing, matter\textquotedblright{}, inherited from Old High German \emph{sahha}, and \emph{Buch} \textquotedblleft{}book\textquotedblright{}, from Old High German \emph{buoh} (EtymDB).
\end{cousinweb}

\begin{grammarlens}[title={What you've built}]
One more word for telling, explaining and describing at B1.
\end{grammarlens}

\subsection*{Your turn}
\begin{itemize}
\raggedright
  \item \emph{Say it:} \emph{das Sachbuch}
  \item \emph{Say it:} \emph{das Sachbuch}, once more
  \item \emph{Say it:} a sentence of your own with \emph{das Sachbuch}
  \item \emph{Recall:} say the German for biography, then the German for short story, then say \emph{das Sachbuch} again
\end{itemize}

\begin{tcolorbox}[breakable,colback=teal!4,colframe=teal!35!black,title={Before you move on}]
What does das Sachbuch mean? (\textquotedblleft{}Non-fiction book\textquotedblright{}.) Say it once more.
\end{tcolorbox}
\section[der Bösewicht]{der Bösewicht — villain}

\label{lesson:GE-C520-boesewicht}



\begin{quote}
Before the new one: say the German for short story, then the German for non-fiction book.
\end{quote}

\subsection*{You'll want to know: der Bösewicht}
\textbf{der Bösewicht} — \textquotedblleft{}villain\textquotedblright{}.

\begin{quote}
Im Film ist der Bösewicht am Ende im Gefängnis.
\end{quote}

\begin{quote}
\emph{In the film the villain ends up in prison.}
\end{quote}

Plural: \emph{die Bösewichte}. Half joking when said of real people.

\begin{cousinweb}
\emph{Bösewicht} is built on \emph{böse} \textquotedblleft{}evil\textquotedblright{}, from Proto-Germanic \emph{bausijaz}, and \emph{Wicht} \textquotedblleft{}creature, wretch\textquotedblright{}, inherited from Old High German \emph{wiht} (EtymDB).
\end{cousinweb}

\begin{grammarlens}[title={What you've built}]
One more word for telling, explaining and describing at B1.
\end{grammarlens}

\subsection*{Your turn}
\begin{itemize}
\raggedright
  \item \emph{Say it:} \emph{der Bösewicht}
  \item \emph{Say it:} \emph{der Bösewicht}, once more
  \item \emph{Say it:} a sentence of your own with \emph{der Bösewicht}
  \item \emph{Recall:} say the German for short story, then the German for non-fiction book, then say \emph{der Bösewicht} again
\end{itemize}

\begin{tcolorbox}[breakable,colback=teal!4,colframe=teal!35!black,title={Before you move on}]
What does der Bösewicht mean? (\textquotedblleft{}Villain\textquotedblright{}.) Say it once more.
\end{tcolorbox}
